\documentclass[a4paper]{article}
\usepackage{luatexja}
\usepackage{amsmath}
\usepackage{amssymb}
\usepackage{amsfonts}
\usepackage{bm}
\usepackage{amsthm}
\usepackage{geometry}
\usepackage{enumitem}
\usepackage{booktabs}
\usepackage{mathtools}
\usepackage{hyperref}

\geometry{margin=25mm}

\newtheorem{theorem}{定理}[section]
\newtheorem{proposition}[theorem]{命題}
\newtheorem{lemma}[theorem]{補題}
\newtheorem{corollary}[theorem]{系}
\theoremstyle{definition}
\newtheorem{definition}[theorem]{定義}
\newtheorem{example}[theorem]{例}
\theoremstyle{remark}
\newtheorem{remark}[theorem]{注意}
\renewcommand{\proofname}{証明}

\newcommand{\Jnorm}{J_{\rm norm}}
\newcommand{\Jbound}{\frac{3\pi^2}{8}}

\title{二重時空四値論理（D4L）：\\
振幅・四状態・幾何干渉・双対ヒルベルト空間}

\author{https://github.com/hypernumbernet}
\date{\today}

\begin{document}
\maketitle

\begin{abstract}
古典二値論理は、すでに決まった世界の論理である。命題は真か偽かであり、矛盾からは何でも従う。二重時空理論は、決まった世界の理論ではない。すべての質量構成は通常と双対の一対のフレームであり、その対の不変量は符号付きの高さ \(\Jnorm\in[-1,1]\) ただ一つである。二重時空四値論理（D4L）は、その高さの論理である。

命題は許容ねじれ配置である。測定は \(\Jnorm\) を読み、四つの名前
\(\lvert T\rangle,\lvert U\rangle,\lvert F\rangle,\lvert B\rangle\)
はその同じ数の符号と飽和であって、外部から足したラベル代数ではない。否定は幾何の通常–双対スワップである。合接と選言は高さの上の \(\min\) と \(\max\) である。配置とそのスワップはともに正の壁に坐れないので、矛盾は爆発しない。

同じ幾何の双対セクターはヒルベルト空間 \(\mathbb{C}^2\) を担う。振幅の演算とヒルベルトの演算は一つの論理に属し、内部辞書で結ばれるが、同一視はしない。キリング形式は不定であり、自己重なりはボルン確率ではなく、四ラベルは \(\mathbb{C}^2\) の PVM ではない。同期と非反駁という二つの指定述語が高さを帰結関係にする。区間の二値制限は、その関係が記録するために作られた状況をホストできない。本稿の主張は機械検証されている。
\end{abstract}

\section{導入}
\label{sec:position}

二値論理は、すでに決まった事実の自然な言語である。命題は真であるか、偽である。両方であれば、他のすべての命題が従う。それが爆発原理であり、古典的帰結が生きた矛盾を許せない理由である。

二重時空理論は、決まった事実を記述しない。すべての質量粒子が担うミンコフスキー・フレームの対——測定する通常フレームと、それに対して時間反転した双対フレーム——を記述する \cite{DST2025,Diophantine2026}。三つの空間軸のそれぞれで、その対は二つのラピディティ \((\alpha_a,\beta_a)\) で符号化される。通常セクターのブーストと、双対セクターの回転である。その対のすでに証明済みの不変量は、符号付きの高さである。二つのフレームが完全に釣り合えば高さは消える。通常セクターが優勢なら高さは正、双対セクターが優勢なら負である。壁では構成は飽和する。すべての軸が純ブーストであるか、すべての軸が純回転である。

これら四つの状況は、代数に後から貼り付けたラベルではない。幾何がすでに供給している量の、符号と飽和である。二重時空四値論理とは、その量を命題の観測量とし、四つの名前を観測量が行うことに任せる、という決定である。

結果の論理は二つの層をもち、混同してはならない。

\emph{振幅層}は、許容配置を命題、正規化高さ \(\Jnorm\) を観測量、\(\min/\max/-\) をスカラー結合子とする。二つの配置の重なりは、PGA コアにすでに存在するパラメータ・キリング形式である。干渉はねじれ双ベクトルの幾何交換子である。代数に何も足さない。論理は、代数がすでに計算しているものを読む。

\emph{双対ヒルベルト層}は双対セクターの運動学である。\(\mathbb{R}^3\) 上のユークリッド対、循環生成子の四元数表、およびその \(\mathbb{C}^2\) 上の表現。閉部分空間の格子は、二次元スピノルのバーコフ--フォン・ノイマン論理である \cite{BvN1936}。それは本物の量子論理である。振幅層の別表記ではない。

内部辞書が二つの層を結ぶ。演算は同一視しない。キリング形式はスピノル内積ではなく、\(\min/\max\) は \(L(\mathbb{C}^2)\) の交わり・和ではなく、四ラベルは射影値測度ではない。\(\hbar\) もなく、\(J\) からのボルン則の導出もない。

本稿はその順でこの絵を展開する。第~\ref{sec:torsion}節で幾何の担体と、四状態を可能にする上限を固定する。第~\ref{sec:states}節で状態に名前を付け、釣り合いラベルを第二の観測量である質量で分け、第五の名前は置かない。第~\ref{sec:conn}節と第~\ref{sec:geo}節でスカラー結合子、二つの順序、幾何演算を与える。第~\ref{sec:hilbert}節で双対ヒルベルト層を取り出し、第一層と同一視しない辞書を述べる。第~\ref{sec:cons}節でどの高さが成り立つかを言い、第~\ref{sec:examples}節で区間の二値制限が D4L の記録する状況をホストできない理由を示す。第~\ref{sec:regime}節で DST の証明成果物とディオファントス・レジームを同じ四つの名前で読む。第~\ref{sec:claims}節で理論の到達と限界を述べる。

\section{命題の担体としてのねじれ}
\label{sec:torsion}

D4L の命題は文ではない。幾何の構成である。六つのラピディティ、同値にねじれ双ベクトルとそれが生成するローター。観測量は、付随する開発ですでに証明済みの単一のスカラーである \cite{Diophantine2026}。

ローター対に対して \(\Omega=R_{\rm usual}^\dagger R_{\rm dual}\) と書き、\(\Omega_{\rm biv}\) を主対数とする。本稿で用いるねじれスカラーのパラメータ形は
\begin{equation}
\label{eq:J-raw}
J = \frac12\sum_{a=1}^3\bigl(\alpha_a^2-\beta_a^2\bigr)
\end{equation}
である。古い草稿に現れる組み合わせ \(\frac1{16}B(\Omega_{\rm biv},\Omega_{\rm biv})\) は、標準的な生成子展開のもとで \eqref{eq:J-raw} と定数倍だけ異なる。ここでの有界性の主張は常に \(J\) とその正規化 \(\Jnorm\) を指す。

任意の六つ組が許されるわけではない。同一軸上の通常ラピディティと双対ラピディティは直角を超えて加算してはならない。それが双対写像の反同期上限である。非負性とあわせて、コンパクトな錐を切り出す。

\begin{definition}[許容構成]
\label{def:admissible}
構成が\emph{許容}であるとは、すべての軸で
\[
\alpha_a\ge 0,\qquad \beta_a\ge 0,\qquad \alpha_a+\beta_a\le \frac{\pi}{2}
\]
が成り立つことである。
\end{definition}

その錐の上で、符号付き高さは無限遠へ走り出せない。二次形式 \(J\) は錐の体積で上から押さえられ、定数倍で上限を \(1\) に直す。

\begin{theorem}[ねじれ上限]
\label{thm:torsion-bound}
定義~\ref{def:admissible} の許容錐の上で
\[
|J|\le\Jbound,\qquad |\Jnorm|\le 1,
\]
ここで
\begin{equation}
\label{eq:Jnorm}
\Jnorm := \frac{8}{3\pi^2}\,J
\end{equation}
である。等号 \(|\Jnorm|=1\) は、すべての軸が純ブーストであるとき、またはすべての軸が純回転であるときに限り成り立つ。許容錐上の \(\Jnorm\) の像は区間 \([-1,1]\) 全体である。主枝条件 \(|\alpha_a+\beta_a|\le\pi/2\) だけを残して負のラピディティを許すと、上限は偽である。
\end{theorem}

したがって D4L の四つの名前は切り詰めではない。\([-1,1]\) のすべての値が達成され、両壁も達成される。ただし、すべての軸が同じ向きに揃う構成によってのみである。混合極端——ある軸が純ブースト、別の軸が純回転——は決して \(|\Jnorm|=1\) に届かない。

コンパクト化 \(N\) に対する離散上限はより鋭い \cite{Diophantine2026}：
\[
|\Jnorm|\le\Bigl(\frac{4\lfloor N/4\rfloor}{N}\Bigr)^2.
\]
\(4\mid N\) のとき右辺は \(1\) であり、両符号が達成される。\(4\nmid N\) のとき、すべての許容格子点は厳密不等式 \(|\Jnorm|<1\) を満たす。正の壁はそこにない。

\emph{振幅}とは、許容連続構成 \(p=(\alpha_a,\beta_a)\) である。そのねじれ双ベクトルとローターは
\[
\Omega(p)=\sum_{a=1}^3\Bigl(\frac{\alpha_a}{2}\,B^+_a+\frac{\beta_a}{2}\,B^-_a\Bigr),
\qquad
R(p)=\exp\Omega(p)
\]
である。測定と崩壊は
\[
\langle p\rangle:=\Jnorm(p),\qquad
\mathrm{collapse}(p)
\in\{\lvert T\rangle,\lvert U\rangle,\lvert F\rangle,\lvert B\rangle\}
\]
である。随伴は通常–双対スワップ \(p^\dagger\) であり、\((\alpha_a)\) と \((\beta_a)\) を入れ替える。スワップは対合であり、観測量の符号を反転する：\((p^\dagger)^\dagger=p\) かつ \(\langle p^\dagger\rangle=-\langle p\rangle\)。すべてのラベルはある振幅によって実現される。

\section{符号付き高さから四状態へ}
\label{sec:states}

\(\Jnorm\) が \([-1,1]\) のすべての値をとることが分かれば、四つの名前はその区間の、符号と飽和による分割である。

\begin{definition}[D4L 状態]
\label{def:states}
\(\Jnorm\in[-1,1]\) の上で
\begin{align*}
\lvert T\rangle &\colon \Jnorm = 0
  &&\text{完全な通常–双対バランス、}\\
\lvert U\rangle &\colon 0 < \Jnorm < 1
  &&\text{ブースト優勢・未飽和、}\\
\lvert F\rangle &\colon \Jnorm = +1
  &&\text{全軸が純ブースト、}\\
\lvert B\rangle &\colon -1\le \Jnorm < 0
  &&\text{回転優勢。最深点 \(\Jnorm=-1\) を含む。}
\end{align*}
\end{definition}

正の極端は False と名付けられた単元である。負の極端は Both の内部に残す。これは論理上の規約であり、錐の欠落ではない。両壁は存在し、一意である。許容配置の上で、\(\lvert F\rangle\) への崩壊は構成が純双曲であることと同値であり、\(\Jnorm=-1\) は純楕円であることと同値である。

内部の \(\lvert B\rangle\) は範囲であり、単一の固定点ではない。その範囲の最深点が壁 \(\Jnorm=-1\) である。したがって四つの名前は随伴のもとで互いの関数ではない。最深の \(\lvert B\rangle\) は \(\lvert F\rangle\) へ、内部の \(\lvert B\rangle\) は \(\lvert U\rangle\) へ行く。ラベルの上の写像としての否定は存在しない。

釣り合いラベル \(\lvert T\rangle\) も、幾何的には一点ではない。符号なしねじれ質量
\[
M(p)=\frac12\sum_{a=1}^3\bigl(\alpha_a^2+\beta_a^2\bigr),\qquad
M_{\mathrm{norm}}=\frac{8}{3\pi^2}\,M
\]
は第二の観測量であり、第五の名前ではない。許容錐上で \(0\le M_{\mathrm{norm}}\le 1\)。通常–双対スワップは \(M\) を保存し、一様スケールは \(M\) を \(k^2\) 倍する。高さは質量を超えない。\(\lvert J\rvert\le M\)、同値に \(\lvert\Jnorm\rvert\le M_{\mathrm{norm}}\)。等号は一方のセクターが完全に消えるときに限り成り立つ。すべての双対ラピディティが零であるか、すべての通常ラピディティが零であるか。質量に対する高さの飽和は、セクターの純度である。

高さ零はなお単一の幾何条件ではない。錐の原点は真空である。すべての軸で \(\alpha_a=\beta_a\) は軸ごとの釣り合い質量である。第三の座席がある。ある軸のブーストが別の軸の回転を打ち消し、軸ごとの等式なしに \(\Jnorm=0\) かつ \(M>0\) となり得る。四状態の分類器は高さだけを見、この三つを区別できない。

真空と釣り合い質量は振幅上の述語である：
\begin{align*}
\mathrm{IsVacuum}(p)
  &\iff \mathrm{collapse}(p)=\lvert T\rangle \wedge M(p)=0,\\
\mathrm{IsBalancedMassive}(p)
  &\iff \mathrm{collapse}(p)=\lvert T\rangle \wedge 0<M(p).
\end{align*}
真空は六つのラピディティがすべて消えることと同値であり、\(\Jnorm=0\) を含意する。逆は偽である。ビールの釣り合い種は軸ごとの座席に坐る。異軸の座席は、同じラベルの別の幾何的理由である。

\section{爆発しない結合子}
\label{sec:conn}

符号付き高さの上のスカラー演算は、実数直線がすでに持っているものである：
\[
\neg j = -j,\qquad
j\wedge k = \min(j,k),\qquad
j\vee k = \max(j,k).
\]
否定は新しい結合子ではない。通常–双対スワップ \((\alpha_a,\beta_a)\mapsto(\beta_a,\alpha_a)\) であり、\(J\mapsto -J\)、したがって \(\Jnorm\mapsto-\Jnorm\) を送る。ド・モルガン双対と分配則は、\(\mathbb{R}\) 上の \(\min/\max\) のものである。

論理的な中身は、高さとそれ自身の否定の合接が\emph{できない}ことにある。

\begin{theorem}[非爆発]
\label{thm:nexplosion}
すべての \(j\in[-1,1]\) に対して
\[
j\wedge\neg j = -|j|\le 0.
\]
ゆえに \(p\wedge\neg p\) は \(\lvert T\rangle\) または \(\lvert B\rangle\) であり、決して \(\lvert F\rangle\) に飽和しない。
\end{theorem}

双対の恒等式は \(j\vee\neg j=\lvert j\rvert\) である。排中は釣り合いでのみ同期し、区間の上ではちょうど二つの壁で反駁される。

配置とその通常–双対スワップは、ともに正の壁に坐れない。矛盾はしたがって回転優勢、または釣り合いの状況である。一意の全ブースト世界ではない。D4L が生きた矛盾をホストしつつ、無関係な原子を \(\lvert F\rangle\) へ強制しないのはそのためである。

区間 \([-1,1]\) は二つの前順序を担い、一つに潰してはならない。\emph{高さ}の順序は実数直線の通常の順序である。\(j\preceq_\tau k\) は \(j\le k\) と同値であり、ブースト優勢と回転優勢を比べる。\emph{情報}の順序は飽和を比べる。\(j\preceq_\iota k\) は \(\lvert j\rvert\le\lvert k\rvert\) と同値である。情報順序では \(j=0\) が唯一の底であり、二つの頂はちょうど壁 \(j=\pm 1\) である。同値に、情報の底はラベル \(\lvert T\rangle\)、二つの頂は \(\lvert F\rangle\) と最深の \(\lvert B\rangle\) である。内部の \(\lvert B\rangle\) は情報順序で最深の \(\lvert B\rangle\) より真に下にある。

したがって四つのラベルはベルナップの \(\{\bot,t,f,\top\}\) と一致しない。双束同型は主張しない。高さと飽和は一つの数の二つの読みであり、独立な二つの格子ではない。

\section{重なりと干渉}
\label{sec:geo}

スカラー結合子は測定のあとで動く。測定の前に、二つの振幅はすでに代数の中で相互作用する。可換な相互作用が重なり、非可換な相互作用が干渉である。
\[
\langle p\mid q\rangle := 8\sum_{a=1}^3\bigl(\alpha_a\alpha'_a-\beta_a\beta'_a\bigr),
\qquad
[\Omega(p),\Omega(q)] := \Omega(p)\Omega(q)-\Omega(q)\Omega(p).
\]
ローター合成 \(R(p)R(q)\) は PGA の積である。ベイカー--キャンベル--ハウスドルフ公式を通して第三のねじれ配置へは戻さない。

重なりは対称であり、自己重なりは観測量の定数倍である。\(\langle p\mid p\rangle=16\,J(p)\)、同値に \(\Jnorm(p)\) の定数倍。異軸のねじれ双ベクトルは可換とは限らない。\(p\) が軸 \(0\) の純ブースト、\(q\) が軸 \(1\) の純回転なら \([\Omega(p),\Omega(q)]\ne 0\)。同軸の双曲生成子と循環生成子は可換なので、対応する一軸干渉は消える。逆は偽である。干渉が消えても単一の共通軸を意味しない。たとえば二軸ブーストは自己交換子が消え、単軸ではない。

自己重なりは \(J\) の定数倍なので、\(J\) の符号を受け継ぐ。純ブースト極端は \(\langle p\mid p\rangle>0\)、純回転極端は \(\langle q\mid q\rangle<0\) を与える。対はしたがって不定である。符号付きの量はボルン確率ではあり得ず、対は第~\ref{sec:spinor}節のスピノル内積でもない。ローターのサンドイッチ評価は重力チャートから独立である。退化二次形式の保存は主張しない。幾何合成は PGA 積として結合的であり、格子演算ではない。

\section{ヒルベルト空間としての双対セクター}
\label{sec:hilbert}

双対セクターはすでにユークリッド三次元空間、四元数表、\(\mathbb{C}^2\) 上のスピノル表現を含んでいる。本節はその運動学を D4L の第二層として読み、第一層との辞書を述べる。空間軸は \(a=0,1,2\) で添字付けする。

\subsection{二つの定符号セクター}
\label{sec:sector}

構成は \(\alpha=0\) のとき双対セクター、\(\beta=0\) のとき通常セクターである。ユークリッド \(\mathbb{R}^3\) の対応する埋め込みを \(p_\beta\)、\(p_\alpha\) と書く。双対セクター上でキリング重なりは
\[
\langle p_\beta\mid p_\gamma\rangle
=-8\langle\beta,\gamma\rangle_{\mathbb{R}^3}
\]
に制限され、したがって負定値である。通常セクター上では \(+8\langle\alpha,\gamma\rangle_{\mathbb{R}^3}\) に制限され、正定値である。内積空間は各セクターの線形包であり、許容立方体 \(\alpha_a,\beta_a\ge 0\)、\(\alpha_a+\beta_a\le\pi/2\) ではない。重ね合わせはその包についての主張である。

\subsection{四元数表}
\label{sec:quat}

双対セクターの循環生成子を \(\Gamma_a\) と書く。それらは \(\Gamma_a^2=-1\) と、四元数代数の非対角積を満たす：
\[
\begin{aligned}
\Gamma_0\Gamma_1&=\Gamma_2,\\
\Gamma_1\Gamma_2&=\Gamma_0,\\
\Gamma_2\Gamma_0&=\Gamma_1,
\end{aligned}
\qquad
\begin{aligned}
\Gamma_1\Gamma_0&=-\Gamma_2,\\
\Gamma_2\Gamma_1&=-\Gamma_0,\\
\Gamma_0\Gamma_2&=-\Gamma_1.
\end{aligned}
\]
恒等式は \(G(3,1,1)\) のベクトル反交換子から従う。純双対構成に対して
\[
\Omega(p_\beta)=\sum_a(\beta_a/2)\,\Gamma_a
\]
であり、軸 \(0\) では \((\theta/2)\,\Gamma_0\) に潰れる。

\subsection{スピノル・ヒルベルト空間}
\label{sec:spinor}

双対スピノル空間は、標準エルミート積をもつ \(\mathbb{C}^2\) である。パウリ行列 \(\sigma_x,\sigma_y,\sigma_z\) は通常のものである。表現 \(\rho(\Gamma_a):=-i\sigma_a\) は、\(IJ=K\) が \(G(3,1,1)\) の循環積と一致し、双対ローターが標準の \(\mathrm{SU}(2)\) 因子 \(\exp(-i\beta\cdot\sigma/2)\) へ送られるように選ぶ。各 \(\rho(\Gamma_a)\) は平方が \(-1\)、循環恒等式 \(\rho(\Gamma_0)\rho(\Gamma_1)=\rho(\Gamma_2)\) が成り立ち、各生成子は反エルミートである。

したがって双対ローター行列は
\[
U(\beta)
=\exp\Bigl(\sum_a(\beta_a/2)(-i\sigma_a)\Bigr)
=\exp(-i\beta\cdot\sigma/2)
\]
である。単一軸ではロドリゲス公式 \(\cos(\theta/2)\,I+\sin(\theta/2)\,(-i\sigma_a)\) である。一般の双対ラピディティに対しても、\(\beta\) が定める単位循環結合に沿って同じ三角形式が成り立つ。

\begin{theorem}[双対ローターは \(\mathrm{SU}(2)\) に属する]
\label{thm:su2}
任意の双対ラピディティ \(\beta\in\mathbb{R}^3\) に対し、行列 \(U(\beta)\) は行列式 \(1\) のユニタリである。ゆえに \(U(\beta)\in\mathrm{SU}(2)\)。純双対構成では \(G(3,1,1)\) のバナッハ指数が同じロドリゲス公式に潰れ、双対セクターの指数と行列ローターはすべての純双対ラピディティで一致する。
\end{theorem}

通常セクターは符号が反対の双対である。相異なる双曲生成子は反交換し、それぞれ平方が \(+1\) なので、結合 \(\Gamma^+(\alpha)=\sum_a\alpha_a B^+_a\) は \(\Gamma^+(\alpha)^2=\lVert\alpha\rVert^2\) を満たす。したがって純通常ラピディティのバナッハ指数は双曲ロドリゲス公式
\[
\exp\bigl(\Omega(p_\alpha)\bigr)
=\cosh(\lVert\alpha\rVert/2)\,I
+\sinh(\lVert\alpha\rVert/2)\,U^+(\alpha)
\]
であり、軸 \(0\) では \(\cosh(\theta/2)\,I+\sinh(\theta/2)\,B^+_0\) である。同じ生成子を \(-i\sigma_a\) ではなくパウリ行列で表すと、通常ローターは
\[
V(\alpha)=\exp(\alpha\cdot\sigma/2)
\]
へ送られる。任意の通常ラピディティに対し \(V(\alpha)\) は行列式 \(1\) のエルミート行列であり、行列指数はすべての純通常構成でバナッハ指数と一致する。非零のブーストはユニタリではない。通常ローターは \(\mathrm{SU}(2)\) の元ではない。行列式 \(1\) のエルミート行列の集合を乗法について群であるとは主張しない。

これが正の壁の運動学である。高さは、一方のセクターが消えるときに限り質量と等しい。通常の指数は \(\lvert F\rangle\) を飽和し、双対の指数は最深の \(\lvert B\rangle\) を飽和する。

混合した通常／双対、または並進生成子に対して、\(G(3,1,1)\) のバナッハ指数をこの行列指数と一般に同一視することは主張しない。単一の純セクターの外でベイカー--キャンベル--ハウスドルフ公式を制御することが必要になるからである。通常ブーストと双対回転の積を \(\mathrm{SL}(2,\mathbb{C})\) の因子分解と同一視もしない。

\(\dim_{\mathbb{C}}\mathbb{C}^2=2\) なので、射影測定の直交する結果は高々二つである。相補的な光線の対は存在する。双対スピノルの座標軸は直交する。四本の互いに直交する光線は存在しない。格子の原子——一次元部分空間——は無限個あり、複素射影直線の各点に一本ある。したがってヒルベルト層は四値論理ではない。四つの D4L ラベルは、それらの原子を数え上げることすらできない。

\subsection[C2 の量子論理]{\(\mathbb{C}^2\) の量子論理}
\label{sec:ql}

量子命題は \(\mathbb{C}^2\) の線形部分空間である。交わりが meet、和が join、直交補空間が否定である。有限次元なのですべての部分空間は閉である。結果の格子はオーソモジュラーである。\(A\le B\) なら \(A\vee(A^\perp\wedge B)=B\)。分配的ではない。\(e_0=(1,0)\)、\(e_1=(0,1)\)、\(d=(1,1)\)、\(\ell_v=\mathbb{C}v\) として
\[
\ell_{e_0}\wedge(\ell_{e_1}\vee\ell_d)
\neq
(\ell_{e_0}\wedge\ell_{e_1})\vee(\ell_{e_0}\wedge\ell_d)
\]
である。左辺は \(\ell_{e_0}\)、右辺は \(\{0\}\)。

第~\ref{sec:conn}節のスカラー結合子は分配的である。分配的な格子は、次元が少なくとも \(2\) のヒルベルト空間の部分空間格子と同型ではあり得ない。\(\mathbb{C}^2\) の互いに直交する一次元部分空間は三つ存在しない。したがって四つの D4L ラベルは四本の直交光線ではあり得ず、振幅結合子はバーコフ--フォン・ノイマン格子を表示できない。

\subsection{射影子とディラック・スピノル}
\label{sec:proj}

代数の原始冪等元は、平方が \(+1\) の生成子から作られなければならない。\(g^2=+1\) なら \((1+g)/2\) は冪等である。軸 \(0\) の双曲生成子を \(j:=e_0e_1\) と書く（したがって \(j^2=+1\)）。この \(j\) は時間ベクトル \(\gamma^0=e_0\) ではない。符号 \((-,+,+,+)\) では後者の平方は \(-1\) である。作業用の射影子
\[
P_{\mathrm{spin}}=\frac{1+e_0e_1}{2},\qquad
P_L=\frac{1-e_1}{2},\qquad
P_R=\frac{1+e_1}{2}
\]
はそれぞれ冪等であり、\(P_L+P_R=1\) である。\(i^2=-1\) の \(\mathrm{Cl}(3,1)\) 擬スカラー \(i\) から作った \((1\pm i)/2\) は冪等ではない。\(\bigl((1-i)/2\bigr)^2=-i/2\)。積 \(P_{\mathrm{spin}}P_R\) も冪等ではない。\(e_0e_1\) と \(e_1\) が反交換するからである。

可換な平方 \(+1\) の対は存在する。\(e_1\) と \(e_0e_2\) である。それらの積射影子
\[
P=\frac{1+e_1}{2}\cdot\frac{1+e_0e_2}{2}
\]
は冪等かつ非零である。それが生成する主左イデアルの実次元はちょうど \(8\) である。それはミンコフスキー部分代数上の四次元左イデアル \(\mathrm{Cl}(3,1)\,P\) と、零生成子 \(e_4\) によるそのコピーとの直和であり、後者は \(G(3,1,1)\) 加群としての真部分加群を与える。したがって八次元イデアルは PGA 加群として可約である。\(8\) という数はディラック・スピノル \(\mathbb{C}^4\) の実次元と一致するが、一致は零方向による倍加であって、八次元イデアルが既約ディラック加群であることではない。ミンコフスキー核の \(\mathrm{Cl}(3,1)\) 既約性、およびいずれかのイデアルと行列表現との同一視は主張しない。

符号 \((-,+,+,+)\) のディラック行列は \(\mathrm{Cl}(3,1)\) の次数 \(1\) 生成子 \(\gamma^\mu\) である。それらは
\[
\{\gamma^\mu,\gamma^\nu\}=2\eta^{\mu\nu},\qquad
\eta=\mathrm{diag}(-1,+1,+1,+1)
\]
を満たす。とくに \((\gamma^0)^2=-1\) なので \(\gamma^0\neq j\)。同じクリフォード関係は、ディラック・スピノル空間 \(\mathbb{C}^4\) に作用する \(M_4(\mathbb{C})\) の行列によって実現される。その空間の複素次元は \(4\) であり、したがって双対スピノル \(\mathbb{C}^2\) ではない。実ベクトル空間としてはワイル分解 \(\mathbb{C}^4\simeq\mathbb{C}^2\oplus\mathbb{C}^2\) がある。分解の代数的な理由はカイラリティである。\(\gamma^5=-i\gamma^0\gamma^1\gamma^2\gamma^3\) とすれば \((\gamma^5)^2=I\) かつ \(\{\gamma^5,\gamma^\mu\}=0\) である。ワイル射影子
\[
P_\pm=\frac{1\pm\gamma^5}{2}
\]
は相補的な冪等元である。上のワイルブロックは \(\gamma^5\) の \(+1\) 固有空間、下のブロックは \(-1\) 固有空間である。その分解は辞書であり、双対セクターの運動学をカイラルなディラック成分と同一視するものではない。

同じクリフォード積は、すでにミンコフスキー波動演算子の一次因子を含む。四元運動量 \(p\) に対し
\[
Q(p)=-p_0^2+p_1^2+p_2^2+p_3^2,\qquad
\gamma(p)=\sum_{\mu}p_\mu\gamma^\mu
\]
と書く。\(\mathrm{Cl}(3,1)\) の関係は代数において \(\gamma(p)^2=Q(p)\) を与え、\(\mathbb{C}^4\) 上の行列としても同じ恒等式が成り立つ。質量殻 \(Q(p)=-m^2\) は、\((\gamma(p)-im)(\gamma(p)+im)\) と因数分解される二次式の消滅である。したがって \(\gamma(p)\Psi=im\Psi\) を満たすスピノルはその殻の上にある。これが運動量空間のディラック方程式であり、独立した要請ではない。ワイル表示では \(\gamma(p)\) はブロック反対角である。質量が零なら二つのカイラルブロックは分離し、質量が非零ならそれらは結合する。ディラック・スピノルをなす二つの \(\mathbb{C}^2\) は、独立した二つの粒子ではなく、一つの質量フェルミオンの二つのカイラル半体である。カイラリティ \(\gamma^5\) は質量 \(m\) の解を質量 \(-m\) の解へ送る。静止系の解は存在する。時空上の偏微分方程式、全ローター群のもとでの共変性、および \(m\) をねじれ遅れと同一視することは主張しない。

\subsection{内部辞書}
\label{sec:dict}

二つの層は幾何を共有する。演算は共有しない。

\begin{center}
\begin{tabular}{@{}ll@{}}
\toprule
振幅層 & 双対ヒルベルト層 \\
\midrule
純双対振幅 & \(\beta\in\mathbb{R}^3\) および \(U(\beta)\in\mathrm{SU}(2)\) \\
純通常振幅 & \(\alpha\in\mathbb{R}^3\) および行列式 \(1\) のエルミート・ブースト \(V(\alpha)\) \\
通常–双対スワップ & 双対セクターを離れる。ヒルベルト随伴ではない \\
キリング重なり & スピノル内積とは別 \\
一軸での干渉消滅 & 可換生成子（必要条件に過ぎない） \\
\(\Jnorm\) の四状態崩壊 & 連続スカラーの粗視化 \\
\(\min/\max\) & \(L(\mathbb{C}^2)\) の交わり・和ではない \\
双対 \(\mathbb{C}^2\) & ディラック \(\mathbb{C}^4\) ではない \\
\bottomrule
\end{tabular}
\end{center}

四つの D4L ラベルに \(\mathbb{C}^2\) の一次元部分空間を割り当て、相異なるラベルを互いに直交させる写像は存在しない。四つの名前は区間の分割であり、量子ビットの射影測定ではない。

\section{どの高さが成り立つか}
\label{sec:cons}

四つのラベルは、どの高さが\emph{成り立つ}かを言わなければ論理にならない。D4L は \(\Jnorm\in[-1,1]\) の上に指定述語を一つではなく二つ置く。

\begin{definition}[指定値]
\label{def:desig}
\[
\mathrm{HoldsT}(j)\iff j=0,
\qquad
\mathrm{HoldsNotF}(j)\iff j<1.
\]
前者は同期（ラベル \(\lvert T\rangle\)）である。後者は非反駁（\(\lvert F\rangle\) 以外のすべてのラベル）である。
\end{definition}

二つの述語は異なる問いに答える。同期は、通常フレームと双対フレームが釣り合っているかを問う。非反駁は、構成が一意の全ブースト世界へ駆動されたかどうかだけを問う。回転優勢の構成は同期ではないが、反駁されてもいない。

指定点は情報順序の底であり、高さの順序の頂ではない。二つの述語は結合子を逆向きに読む。合接の非反駁は、非反駁の選言である。交わりが正の壁に届くのは、両方の高さがすでにそこに坐るときだけである。選言の非反駁は、非反駁の合接である。和が壁に届くのは、どちらか一方が届いたときである。同期はより厳しく、かつ吸収する。釣り合いと任意の非負の高さ——正の壁を含む——との合接は、再び釣り合いである。釣り合いと任意の非正の高さとの選言も、再び釣り合いである。とくに \(0\wedge 1=0\) であり、正の壁は合接を吸収しない。

論理式は原子、\(\neg\)、\(\wedge\)、\(\vee\) から生成し、高さの上ですでに固定したスカラー演算で評価する。振幅の高さに含意結合子はない。格子の頂を \(1\) とする \(\min\) のゲーデル剰余は、この論理には合わない。ブースト優勢の高さが同期を含意すると、その剰余は同期そのものを指定値にし、回転優勢の高さが同期を含意すると正の壁 \(\lvert F\rangle\) へ送る。指定値が釣り合いと、その壁の外にある論理は、その剰余を使えない。含意は離散の証明ステータス層（第~\ref{sec:regime}節）にある。

付値は各原子に区間 \([-1,1]\) の高さを与える。帰結は二つの述語に従う。\(\Gamma\models_T\psi\) は、\(\Gamma\) を同期するすべての付値が \(\psi\) を同期することであり、\(\Gamma\models_{\lnot F}\psi\) は同じ条項を非反駁で書いたものである。二つの関係は同じではない。\(P\wedge\neg P\) はあらゆる付値のもとで非反駁なので、空集合はそれを非反駁する。同期するのは \(P\) が同期するときに限るので、空集合はそれを \(T\)-含意しない。原子の同期は \(P\vee\neg P\) の同期を強制する。非反駁は強制しない。負の壁は反駁されていないが、その壁とスワップとの選言は正の壁である。

区間の二値\emph{制限}は一意ではない。D4L の名前 \(\lvert T\rangle,\lvert F\rangle\) は古典の真／偽ではないからである。\emph{名前}フラグメントは \(\{0,1\}=\{\lvert T\rangle,\lvert F\rangle\}\) である。\emph{壁}フラグメントは \(\{-1,+1\}=\{\text{最深の }\lvert B\rangle,\lvert F\rangle\}\) である。否定のもとでの振る舞いは異なる。名前フラグメントでは \(\neg 0=0\) である。釣り合い点は否定の固定点であり、\(\neg 1=-1\) は \(\{0,1\}\) を離れる。したがって名前フラグメントはブールではなく、閉じてすらいない。壁の上では \(\neg\) が二点を入れ替え、固定点はない。壁フラグメントが忠実な古典の核である。D4L を「真と偽」と比べるときに名前 \(\lvert T\rangle\) と \(\lvert F\rangle\) へ切ると、すでに古典論理を離れている。

\section{二値論理がホストできないもの}
\label{sec:examples}

D4L と二値論理の差は、D4L がより強い証明器であることでも、より多くの算術を決めることでもない。二値代数がホストできない充足・帰結・格子の状況を D4L がホストし、その差を対として固定できることである。フラグメント側の空性（または空虚、単元への潰れ）と、D4L 側の実現。

これらは意味論の変更であり、ペアノ算術の新しい定理ではない。無条件の FLT / ビールも、ゲーデル定理の反証も主張しない。

\subsection{否定の固定点}
\label{sec:ex-fp}

古典では \(P\leftrightarrow\neg P\) にモデルがない。壁フラグメントでも同じである。\(j\in\{-1,+1\}\) なら \(j\neq -j\)。D4L では条件 \(j=\neg j\) は \(j=0\) と同値である。一意の解はラベル \(\lvert T\rangle\)、完全な通常–双対バランスであり、定数付値でも振幅でも実現される。

この等式はゲーデル文ではない。自己言及的な\emph{等式} \(P=\neg P\) であり、否定に固定点のない代数はホストできず、D4L はそれを同期に一意に置く。

\subsection{爆発しない矛盾}
\label{sec:ex-efq}

古典二値論理には \emph{ex falso quodlibet} がある。\(\{P,\neg P\}\) から任意の \(Q\) が従う。\(\mathrm{Bool}\) の上では前提がすでに不可能である。\(P\) と \(\neg P\) はともに \(\mathtt{true}\) になれない。壁の上では \(j\) と \(\neg j\) がともに \(+1\) にはなれない。したがって二値の爆発は空虚である。前提を同時に指定できない。

D4L ではできる。\(P\) を原子 \(0\)、\(Q\) を原子 \(1\) とすると
\[
\{P,\neg P\}\not\models_T Q,\qquad
\{P,\neg P\}\not\models_{\lnot F} Q
\]
である。一つの証人は \(P\) を \(0\)、\(Q\) を \(1\) へ送る。矛盾対は同期し、無関係原子は False である。第二の証人は \(P\) を \(-\tfrac12\)（内部の \(\lvert B\rangle\)）、\(Q\) を \(1\) へ送る。対は非反駁であり、無関係原子はやはり False である。連言 \(j\wedge\neg j\) は常に非反駁である。矛盾した観測の対は、無関係な原子を正の壁へ強制しない。

\subsection{\texorpdfstring{制約 \(\Jnorm<1\)}{制約 Jnorm < 1}}
\label{sec:ex-notF}

一意の全ブースト世界に坐ることを禁じられた文は、\(\mathrm{HoldsNotF}\)、すなわち \(\Jnorm<1\) で制約される。名前フラグメントの上ではこの制約は \(\mathrm{HoldsT}\) に潰れる。残る点は \(0\) だけである。D4L では同じ制約が三つのラベル \(\lvert T\rangle\)、\(\lvert U\rangle\)、\(\lvert B\rangle\) によって実現される。非反駁は同期ではない。名前フラグメントはその差を見ない。

\subsection{相補性}
\label{sec:ex-comp}

\(\mathbb{C}^2\) の部分空間格子は分配的でなく、したがってブール代数ではない。互いに直交する原子は高々二つである。四つの D4L ラベルは PVM ではない。分配的なスカラー結合子はその格子を表示できない。双対ヒルベルト層の相補的な光線にブール値を割り当てることはできず、振幅層の四つの名前と同一視することもできない。二つの層は二つの層のままである。

\section{証明ステータスとディオファントス・レジーム}
\label{sec:regime}

DST の証明成果物——核の補題、未閉鎖の bridge、棄却済み診断、帳簿の恒等式——はねじれ配置ではない。四つの名前を再利用する離散ステータスを担うが、観測量は \(\Jnorm\) ではない。この担体の上では、否定は四ラベルの関数であり、含意は表であって \(\min\) の剰余ではない。

否定と格子演算は振幅の代表 \(\lvert T\rangle\mapsto 0\)、\(\lvert U\rangle\mapsto 1/2\)、\(\lvert F\rangle\mapsto 1\)、\(\lvert B\rangle\mapsto -1/2\) に従う。とくに \(F\wedge\neg F=B\)。第二の演算である確立の交わりは、二つの成果物を \(F<B<U<T\) に沿って束ねる。証明済みの核と未閉鎖の bridge の対はまだ未閉鎖である。\(T\) と \(U\) の交わりは \(U\)。含意は
\[
\begin{aligned}
T\to\phi&=\phi,\\
F\to\phi&=T,\\
U\to T&=U,\\
B\to T&=U
\end{aligned}
\]
である。最後の二条が要点である。開いた、または衝突した計画は定理を確立しない。それらは \(\min\) の剰余を拒んだ理由でもある。振幅の代表の上では \(U\to T\) は同期に、\(B\to T\) は正の壁に評価される。表はどちらも非指定のステータスに留める。ベルナップ FOUR 同型は主張しない。ステータス上の指定述語は同じ二つの名前を再利用する。\(\mathrm{HoldsT}_R\) はステータス \(T\)、\(\mathrm{HoldsNotF}_R\) はステータス \(\ne F\) である。

これが、DST が証明したものと証明していないものを記録する代数である。名前の二値論理 \(\{T,F\}\) は未閉鎖の状態 \(U\) をホストできない。壁フラグメント \(\{B,F\}\) は証明済みの核 \(T\) をホストできない。D4L は核 \(T\)、診断 \(F\)、本命 \(U\)、予想 \(U\) の割り当てを実現する。核だけでは予想を \(T\)-含意しない。窓の \(T\) と ill-posed な NoGo の \(B\) は、予想を \(F\) の外へ強制しない。含意 \(U\to T\) は指定値ではない。

同じ代数がビールの残件地図をホストする。閉じた切片（フェルマーおよびダルモン--メレルの位置、符号 \((n,n,5)\) を含む）は \(T\) に坐り、棄却済み診断は \(F\)、帳簿の恒等式は \(B\)、未閉鎖の残件と古典予想は \(U\) に坐る。閉じた切片は古典ビールを \(T\)-含意しない。立方および五乗符号の偶置換は未閉鎖のまま残る。偶数の共通指数は ダルモン--メレル 形にも \( (n,n,5) \) 形にも書き直せず、それらの三つ組は閉じた切片ではなく偶二一致の差残件に属する。連続する完全冪に関する ミハイレスク の定理は、係数条件を閉じても、それを宿す指数形状までは閉じない。ステータス盤は残件の算術的組立を符号化しない。名前付き成果物をすべて \(T\) に強制しても予想を \(T\)-含意せず、未閉鎖残件を一つ \(U\) から \(T\) へ昇格させても予想を主張できない。指数形状の分類器は \(\{T,U,F\}\) に着地し、帳簿の \(B\) を記録しない。双対軸のフェルマー地図も同様に、加法核、実 \(L^p\) 二分法、および \(N_1\) の楕円的有限サンドイッチを \(T\)、格下げされた単軸モジュラーおよび釣り合い種の診断を \(F\)、混合モーター残件と古典 FLT を \(U\) に置く。閉じたフェルマー切片は古典 FLT を \(T\)-含意しない。

DST が単一の上限や棄却済み診断を七つの古典予想の含意として扱ってはならないのは、このためである。レジーム層はその禁止を記録する。未証明の bridge を証明せず、無条件の FLT、ビール、abc も主張しない。

加法モーターと乗法振幅は、同じラベル \(\lvert T\rangle\) の別の座席を占める。純並進はねじれが消え、真空である。釣り合い質量構成は真空ではない。因子が可換でないとき、混合積 \(R\cdot T\) は \(\exp(\Omega_{\mathrm{biv}})\) と同一視しない。釣り合い残件クラスは窓の種を除外する。半窓構成は釣り合い残件ではない。

有限トーラスの上で、同じ四つの名前は離散の読みを得る。\(4\nmid N\) のとき鋭い離散上限は厳密であり、許容点は \(\lvert F\rangle\) に飽和しない。\(4\mid N\) のとき両壁が達成される。全ブーストの格子点が \(\lvert F\rangle\) を実現し、全回転の格子点が最深の \(\lvert B\rangle\) を実現する。純ブーストのスケールは \(\Jnorm\) を \(k^2\) 倍する。同期は固定点である。ラベル \(\lvert U\rangle\) は \(\lvert F\rangle\) へ駆動され得るし、許容錐の外へ出て、崩壊が未定義になり得る。それが共有 no-go の論理的読みであり、新しい算術ではない。

真空はすべての許容スケールの固定点である。小さな釣り合い質量の種は質量 \(k^2 M\) のまま \(\lvert T\rangle\) に留まり、許容であり続け得る。大きな釣り合い種は錐を出る。その錐の退出は符号付き高さの分類器には見えない。\(\Jnorm\) は \(0\) のままなので、四状態ラベルはなお \(\lvert T\rangle\) である。それを見るのは質量と錐述語である。ビールの釣り合い種 \(\theta=\log 2/m\) が錐の内部で \(m\) 倍スケールを生き延び得るのは、このためである。

高さとモジュラー巻数は二つの測定である。実スケールの許容性と非零巻数は同時には成り立たない。その両立不可能性は観測量の分離であり、一つの観測量の上の Both ステータスではない。

\section{自己言及の読み}
\label{sec:godel}

\(G\) をゲーデル文とする。「\(G\) はこの体系で証明可能ではない」。D4L では、証明可能性を同期 \(\Jnorm\to 0\)、すなわち \(\lvert T\rangle\) への駆動として読む。一意の全ブースト世界に坐ることを禁じられた文に対する誠実な制約は、第~\ref{sec:ex-notF}節の \(\mathrm{HoldsNotF}\) そのものである：
\[
\Jnorm(G)<1.
\]
この制約は \(G\) をラベルとして \(\lvert T\rangle\)、\(\lvert U\rangle\)、または \(\lvert B\rangle\) に残す。さらに同期を要求するなら \(\lvert T\rangle\) である。連言は爆発しない（第~\ref{sec:ex-efq}節）。

これは意味論の変更であり、ゲーデルの定理が古典算術において失敗することの証明ではない。負トーションセクター上の完全性の話も、許容ローターの読みであって、一次ペアノ算術の内部でのヒルベルト計画の完了ではない。第~\ref{sec:examples}節の内容は固定されており、本節の読みは新しい不完全性定理ではない。

\section{到達範囲}
\label{sec:claims}

D4L は許容ねじれ配置を命題の担体とし、\(\Jnorm\) を観測量とする。\([-1,1]\) を四つの名前付き状態へ分割し、許容錐上で両壁が達成される。否定を幾何の通常–双対スワップにする。高さとそれ自身の否定の合接は \(\lvert F\rangle\) の外に留まる。釣り合いとの合接は釣り合いへ戻り、非反駁は交わりと和を入れ替える。高さと飽和は一つの区間の上の二つの前順序である。高さは質量を超えず、等号は一方のセクターが消えるときに限る。キリング対は不定であり、自己重なりはボルン確率ではない。双対セクターはユークリッド三次元空間、四元数表、および \(\mathbb{C}^2\) 上の \(\mathrm{SU}(2)\) ローターを与え、その閉部分空間格子はオーソモジュラーかつ非分配的であり、原子は無限個、相補対はあるが四つではない。通常セクターはその双対である。双曲ロドリゲスと、非零ブーストではユニタリでない行列式 \(1\) のエルミート行列ローター。カイラリティ \(\gamma^5\) は平方が単位、ディラック行列と反交換し、相補的なワイル射影子で \(\mathbb{C}^4\) を分ける。ミンコフスキー二次形式は \(\mathrm{Cl}(3,1)\) ですでに平方なので、\(\gamma(p)\Psi=im\Psi\) を満たす運動量空間のスピノルは質量殻 \(Q(p)=-m^2\) の上にある。質量はワイルブロックを結合し、\(\gamma^5\) は \(m\) の符号を反転する。

二つの指定述語と二つの二値フラグメントが、その絵を帰結関係にする。その関係は \(\lvert T\rangle\) に一意の否定固定点をホストし、\(\{P,\neg P\}\) の爆発を拒み、\(\Jnorm<1\) を三ラベルで実現し、DST の三層——証明済みの核、棄却済み診断、未閉鎖の bridge——を記録して、古典予想を \(T\)-含意しない。含意はステータス代数にあり、\(U\to T=U\) である。\(\min\) の剰余はその条項を同期として指定し、\(B\to T\) を \(\lvert F\rangle\) へ送ってしまうからである。ラベル \(\lvert T\rangle\) は幾何的に一点ではない。真空、軸ごとの釣り合い質量、異軸の打ち消しがそこに坐り、質量と軸ごとの等式で分かれ、第五の名前はない。\(4\nmid N\) のとき離散の \(\lvert F\rangle\) は禁じられ、\(4\mid N\) のとき両壁が達成される。増幅はラベル上の部分写像であり、高さ分類器が見えない釣り合い質量の錐退出を含む。

理論は振幅演算とヒルベルト演算を同一視せず、四ラベルを PVM と同一視せず、干渉消滅を共通単軸と同一視せず、生の \(J\) を \([-1,1]\) の値と同一視しない。ボルン則も時空上のディラック偏微分方程式も測定問題の解決も導かない。運動量空間のディラック方程式は第~\ref{sec:proj}節に記したミンコフスキー二次形式の一次因子であり、質量をねじれから定めることではない。可換左イデアルが \(G(3,1,1)\) 加群として零方向で可約であることは第~\ref{sec:proj}節に記した。ミンコフスキー核の \(\mathrm{Cl}(3,1)\) 既約性も、いずれかのイデアルと行列ディラック加群との同一視も主張せず、単一の純セクターを超えてバナッハ指数を行列指数と同一視せず、混合した通常／双対ローターを \(\mathrm{SL}(2,\mathbb{C})\) で因子分解もしない。古典算術におけるゲーデル定理を反証せず、無条件のビール、FLT、abc も主張しない。

以上の主張は機械検証されている。

\begin{thebibliography}{9}
\raggedright

\bibitem{DST2025}
Dual Spacetime Theory: Gravity as Torsion between Dual Spacetimes (2025),
arXiv:2512.xxxxx.

\bibitem{Diophantine2026}
Discrete Biquaternionic Double Spacetime: Rigorous Proofs of Diophantine
Conjectures via Torsional Mismatch and PGA Embedding,
\texttt{dst-diophantine.tex} (2026). Lean companion:
\url{https://github.com/hypernumbernet/DstDiophantine}.

\bibitem{BvN1936}
G.~Birkhoff and J.~von Neumann, The logic of quantum mechanics,
Ann.\ Math.\ 37 (1936).

\bibitem{Godel1931}
K.~G\"odel, \"Uber formal unentscheidbare S\"atze der Principia Mathematica
und verwandter Systeme~I, Monatshefte Math.\ Phys.\ 38 (1931).

\bibitem{Priest2006}
G.~Priest, \emph{In Contradiction: A Study of the Transconsistent},
Oxford University Press (2006).

\end{thebibliography}

\end{document}
